% Mechanically extracted after \mainmatter from the frozen standalone master.

\markboth{Exposé XXV. The Existence Theorem}{Exposé XXV. The Existence Theorem}
% BEGIN inputs/XXV/tex/components/00_title_introduction_and_section1.tex
\chapter*{Exposé XXV\\The Existence Theorem}
\addcontentsline{toc}{chapter}{Exposé XXV. The Existence Theorem}
\begin{center}
\textit{by M. Demazure}
\end{center}

\footnote{Editors' note: Version of 13 October 2024.}
For completeness, in this Exposé we give a proof of the existence theorem for split
groups. Like Chevalley's original proof (``Sur certains schémas de groupes
semi-simples'', \textit{Séminaire Bourbaki}, May 1961, no.~219), it rests on the
existence of complex semisimple algebraic groups of every possible type. Cartier
gave the principle of a more satisfactory proof, establishing directly the
existence of a simply connected split semisimple \(\mathbb Z\)-group associated
with a given Cartan matrix (unpublished).\footnote{Editors' note: Such a principle
was also sketched by B.~Kostant \emph{[Ko66]}; a complete proof, using quantum
groups, was given more recently by G.~Lusztig \emph{[Lu09]}.}
Let us point out, however, that the difficulty is not to give an explicit
construction of a group scheme, but to verify that the group so constructed
really satisfies the required conditions---essentially, that its fibers are
indeed smooth and reductive.

\SGASection{1}{Statement of the theorem}
\SGARefTarget{sga3:xxv:section:1}

\begin{sgastatement}{Theorem 1.1}
Let \(S\) be a nonempty scheme. The functor
\[
  G\longmapsto R(G)
\]
is an equivalence from the category of pinned reductive \(S\)-groups to the
category of pinned reduced root data.
\end{sgastatement}
\SGARefTarget{sga3:xxv:theorem:1:1}

By the uniqueness theorem (Exposé XXIII, 4.1), the preceding statement is
equivalent to the following.

\begin{sgastatement}{\SGAInlineRefTarget{sga3:xxv:corollary:1-2}Corollary 1.2 (Existence of split groups)}
For every reduced root datum \(R\), there exists a split reductive
\(\mathbb Z\)-group \(G\) such that \(R(G)\simeq R\).\footnote{Editors' note:
The work \emph{[BT84]} of F.~Bruhat and J.~Tits also contains a variant of
Chevalley's construction (\emph{loc.\ cit.}, 2.2.3--2.2.5 and \S3.2), which
in particular gives a smooth affine \(\mathbb Z\)-group \(G\) with connected
fibers, possessing a split maximal torus, and whose generic fiber is a reductive
\(\mathbb Q\)-group of type \(R\). The reductivity of the geometric fibers of
\(G\) follows from the study of the unipotent radical of a special fiber in
\emph{loc.\ cit.}, 4.6.12 and 4.6.15 (valid for more general \(G\), associated
with a valued root datum); but it is simpler to deduce it from the description
of \(\mathfrak g=\operatorname{Lie}(G)\), from Exposé XIX, 1.12(iii), and from
the existence of the elements \(w_\alpha(X)\) of \SGARefLink{sga3:xx:3.1:item:iv}{Exposé XX, 3.1(iv)}; see
\emph{[BT84]}, 3.2.1.}
\end{sgastatement}
\SGARefTarget{sga3:xxv:corollary:1:2}

In particular:

\begin{sgastatement}{Corollary 1.3}
Let \(k\) be a field. For every splittable reductive \(k\)-group \(G_k\), there
exists a reductive \(\mathbb Z\)-group \(G\) such that
\[
  G\otimes_{\mathbb Z}k\simeq G_k.
\]
\end{sgastatement}
\SGARefTarget{sga3:xxv:corollary:1:3}

Conversely, first observe that to prove \SGARefLink{sga3:xxv:corollary:1-2}{1.2} it is enough, by \SGARefLink{sga3:xxii:prop:4.3.1}{Exposé XXII, 4.3.1},
and \SGARefLink{sga3:XXI:6.5.10}{Exposé XXI, 6.5.10}, to consider the case in which the root datum \(R\) is
simply connected (and even irreducible, if desired, by \SGARefLink{sga3:XXI:7.1.6}{Exposé XXI, 7.1.6}).

On the other hand, under the hypotheses of \SGARefLink{sga3:xxv:corollary:1:3}{1.3} the group scheme \(G\) has
constant type (since \(\operatorname{Spec}(\mathbb Z)\) is connected), hence type
\(R(G_k)\). By \SGARefLink{sga3:XXIII:5.9}{Exposé XXIII, 5.9}, the validity of \SGARefLink{sga3:xxv:corollary:1:3}{1.3} for a given group \(G_k\)
therefore entails the existence of a split \(\mathbb Z\)-group of type
\(R(G_k)\).\footnote{Editors' note: In fact, \SGARefLink{sga3:XXIII:5.9}{Exposé XXIII, 5.9} is unnecessary:
the present Exposé constructs, for every semisimple \(\mathbb Q\)-group
\(G_{\mathbb Q}\), a semisimple \(\mathbb Z\)-group \(G\) of the same type as
\(G_{\mathbb Q}\), equipped with a split maximal torus \(T\); hence \(G\) is
split by \SGARefLink{sga3:xxii:prop:2.2}{Exposé XXII, 2.2}.}

Thus, to prove \SGARefLink{sga3:xxv:corollary:1-2}{1.2} and hence \SGARefLink{sga3:xxv:theorem:1:1}{1.1}, it is enough to prove \SGARefLink{sga3:xxv:corollary:1:3}{1.3} when \(k\) has
characteristic zero (for example, \(k=\mathbb C\)) and \(G_k\) is simply
connected (and in particular semisimple), together with the following
proposition.

\begin{sgastatement}{Proposition 1.4}
For every simply connected reduced root datum \(R\), there exists a semisimple
algebraic \(\mathbb C\)-group of type \(R\).
\end{sgastatement}
\SGARefTarget{sga3:xxv:proposition:1:4}

One may prove \SGARefLink{sga3:xxv:proposition:1:4}{1.4} as follows. First, there exists a complex semisimple Lie
algebra \(\mathfrak g\) of type \(R\); see, for example, N.~Jacobson,
\textit{Lie Algebras}, Chapter VII, Theorem 5.\footnote{Editors' note:
Let \(\Delta\) be a base of \(R\), and let \(\widetilde{\mathfrak g}\) be the
complex Lie algebra generated by
\((e_\alpha,f_\alpha,h_\alpha)_{\alpha\in\Delta}\), subject to the relations
\([h_\alpha,h_\beta]=0\), \([e_\alpha,f_\beta]=h_\alpha\) if
\(\beta=\alpha\) and \(=0\) otherwise,
\([h_\alpha,e_\beta]=(\alpha^\ast,\beta)e_\beta\), and
\([h_\alpha,f_\beta]=-(\alpha^\ast,\beta)f_\beta\).
In \emph{loc.\ cit.}, \(\mathfrak g\) is defined as the quotient of
\(\widetilde{\mathfrak g}\) by the intersection of the kernels of its
finite-dimensional irreducible representations. In \emph{[Se66]},
\S VI.5, Theorem 9 (see also \emph{[BLie]}, VIII, \S4.3, \SGARefLink{sga3:xxv:section:1}{Theorem 1}), it is
shown that \(\mathfrak g\) is the quotient of \(\widetilde{\mathfrak g}\) by
the relations
\[
 \operatorname{ad}(e_\alpha)^{1-(\alpha^\ast,\beta)}(e_\beta)=0,\qquad
 \operatorname{ad}(f_\alpha)^{1-(\alpha^\ast,\beta)}(f_\beta)=0.
\]
For an explicit description of the structure constants, and in particular
the choice of signs, see \SGARefLink{sga3:XXIII:6.5}{Exposé XXIII, 6.5} and 6.7, as well as
\emph{[Ti66]}, \S4, \SGARefLink{sga3:xxv:section:1}{Theorem 1}, and, for types \(A,D,E\),
\emph{[Sp98]}, 10.2.5.}
Then \(G=\operatorname{Aut}(\mathfrak g)^0\) is a semisimple algebraic
\(\mathbb C\)-group of type \(\operatorname{ad}(R)\).\footnote{Editors' note:
One may suppose \(R\) irreducible, hence \(\mathfrak g\) simple. By a general
argument, \(\operatorname{Lie}(G)\) is the complex Lie algebra of derivations
of \(\mathfrak g\); see \emph{[DG70]}, \S II.4, Proposition 2.3. These
derivations are all inner; see \emph{[BLie]}, I, \S6.1, Corollary 3 to
\SGARefLink{sga3:xxv:section:1}{Proposition 1}. Thus \(\operatorname{Lie}(G)=\mathfrak g\), so \(G\) has no
invariant subgroup of positive dimension and is semisimple. Its root system
is consequently the same as that of \(\mathfrak g\). Moreover, the center of
\(G\) acts both trivially and faithfully on \(\mathfrak g\), and is therefore
trivial; hence \(G\) is adjoint.}
By \textit{Bible}, \S23.1, \SGARefLink{sga3:xxv:section:1}{Proposition 1}, it follows that there exists a
semisimple \(\mathbb C\)-group of type \(R\).

The remainder of this Exposé is devoted to the proof of \SGARefLink{sga3:xxv:corollary:1:3}{1.3} when \(k\) has
characteristic zero and \(G_k\) is semisimple. The proof proceeds in two
stages: the construction of a ``piece of \(\mathbb Z\)-group scheme''
(\S2), and the study of the group obtained by applying ``Weil's theorem''
(\S3).

To avoid confusion, in \S2 we shall not use the abbreviated notation \(X_k\)
for the \(k\)-scheme \(X\otimes_{\mathbb Z}k\), where \(X\) is a
\(\mathbb Z\)-scheme.


% END inputs/XXV/tex/components/00_title_introduction_and_section1.tex

% BEGIN inputs/XXV/tex/components/01_section2_1_through_2_5.tex
\SGASection{2}{Existence theorem: construction of a piece of group}
\SGARefTarget{sga3:xxv:section:2}

\SGANextSectionTarget{sga3:xxv:subsection:2-1}\SGASubsection{2.1}{}
\SGARefTarget{sga3:xxv:paragraph:2:1}

Choose once and for all a splitting of \(G_k\), denoted
\[
  (G_k,T_k,M,R)
\]
(see \SGARefLink{sga3:xxii:def:1.13}{Exposé XXII, 1.13}), a system \(\Delta\) of simple roots of \(R\)
(defining the system \(R^+\) of positive roots), and a Chevalley system
\((X_{\alpha,k})_{\alpha\in R}\) of \(G_k\) (\SGARefLink{sga3:XXIII:6.1}{Exposé XXIII, 6.1} and \SGARefLink{sga3:XXIII:6.2}{6.2})
satisfying the additional condition (see \SGARefLink{sga3:xx:2.6}{Exposé XX, 2.6})
\[
  X_{\alpha,k}X_{-\alpha,k}=1\qquad(\alpha\in R).
\]
Finally, choose on the subgroup of \(M\) generated by \(R\) a total order
compatible with the group structure and such that the roots \(>0\) are the
elements of \(R^+\). We then write the roots as
\[
 -\alpha_n<-\alpha_{n-1}<\cdots<-\alpha_1
 <\alpha_1<\alpha_2<\cdots<\alpha_n.
\]
For \(\alpha\in R\), let \(U_{\alpha,k}\) denote the vector group
corresponding to the root \(\alpha\), and let
\[
 p_{\alpha,k}\colon\mathbb G_{a,k}\xrightarrow{\ \sim\ }U_{\alpha,k}
\]
be the isomorphism of vector groups defined by \(X_{\alpha,k}\).

\SGANextSectionTarget{sga3:xxv:subsection:2-2}\SGASubsection{2.2}{}
\SGARefTarget{sga3:xxv:paragraph:2:2}

The splitting of \(G_k\) includes, in particular, an isomorphism of
\(k\)-groups
\[
 T_k\simeq D_k(M).
\]
Put \(T=D(M)\). This is a \(\mathbb Z\)-torus, and the preceding isomorphism
may be regarded as an isomorphism
\[
 T_k\simeq T\otimes_{\mathbb Z}k.
\]
We have
\[
 \operatorname{Hom}_{\mathbb Z\text{-gr.}}(T,\mathbb G_m)=M,
\]
and shall regard the elements of \(R\subset M\) as characters of \(T\).
Likewise, the elements of \(R^\ast\) will be regarded as morphisms of
\(\mathbb Z\)-groups \(\mathbb G_m\to T\).

\SGANextSectionTarget{sga3:xxv:subsection:2-3}\SGASubsection{2.3}{}
\SGARefTarget{sga3:xxv:paragraph:2:3}

For each \(\alpha\in R^+\), let \(\mathbb G_a(\alpha)\) be a copy of
\(\mathbb G_a\), and consider the \(\mathbb Z\)-scheme
\[
 U=\mathbb G_a(\alpha_1)\times\cdots\times\mathbb G_a(\alpha_n).
\]
If \(U_k\) denotes the unipotent part of the Borel subgroup \(B_k\) of
\(G_k\) defined by \(R^+\), let
\[
 a\colon U\otimes_{\mathbb Z}k\xrightarrow{\ \sim\ }U_k
\]
be the isomorphism of \(k\)-schemes defined by
\[
 a(x_1,\ldots,x_n)
 =p_{\alpha_1,k}(x_1)\cdots p_{\alpha_n,k}(x_n).
\]

\SGANextSectionTarget{sga3:xxv:subsection:2-4}\SGASubsection{2.4}{}
\SGARefTarget{sga3:xxv:paragraph:2:4}

The group law of \(U_k\) is expressed by relations of the form
\[
 a(x_1,\ldots,x_n)a(y_1,\ldots,y_n)=a(z_1,\ldots,z_n),
\]
where, for \(h=1,\ldots,n\), each \(z_h\) is a polynomial
\[
 z_h=x_h+y_h+
 Q_h(x_1,\ldots,x_{h-1},y_1,\ldots,y_{h-1}),
\]
whose coefficients are integers (\SGARefLink{sga3:xxii:lem:5.5.8}{Exposé XXII, 5.5.8}, and \SGARefLink{sga3:XXIII:6.4}{Exposé XXIII,
6.4}). Moreover,
\[
 Q_h(x_1,\ldots,x_{h-1},0,\ldots,0)=0.
\]

Equip \(U\) with the composition law defined by these formulas, which are
indeed ``defined over \(\mathbb Z\)''. Since this law induces the group law
on \(U_k\), it is associative, and \((0)\) is a unit element. Indeed, these
two assertions are expressed by relations among the polynomials \(Q_h\),
and \(\mathbb Z\to k\) is injective. This is a group law: if \((x_i)\) is
a section of \(U\) over an arbitrary scheme \(S\), its inverse \((y_i)\)
is computed recursively by
\[
 y_i=-x_i-Q_i(x_1,\ldots,x_{i-1},y_1,\ldots,y_{i-1}),
\]
formulas which are again ``defined over \(\mathbb Z\)''.

In summary, we have constructed a group law on \(U\) for which the
isomorphism \(a\) above is an isomorphism of groups.

For each \(\alpha\in R^+\), consider the morphism
\[
 p_\alpha\colon\mathbb G_a\longrightarrow U
\]
defined by \(p_\alpha(x)=(x_i)\), where
\[
 x_i=
 \begin{cases}
 x,&\alpha_i=\alpha,\\
 0,&\alpha_i\ne\alpha.
 \end{cases}
\]
It is a closed immersion and a group homomorphism; denote its image by
\(U_\alpha\). Since
\[
 (x_i)=p_{\alpha_1}(x_1)\cdots p_{\alpha_n}(x_n),
\]
\(U\) is identified with the product
\[
 U=U_{\alpha_1}U_{\alpha_2}\cdots U_{\alpha_n}.
\]

\SGANextSectionTarget{sga3:xxv:subsection:2-5}\SGASubsection{2.5}{}
\SGARefTarget{sga3:xxv:paragraph:2:5}

Let \(T=D(M)\) act on each \(U_\alpha\) through the character \(\alpha\).
One verifies immediately that this defines an action of \(T\) on the
group \(U\), and one may form the semidirect product \(B=T\ltimes U\).
There is a canonical isomorphism of \(k\)-groups
\[
 B\otimes_{\mathbb Z}k\xrightarrow{\ \sim\ }B_k.
\]

If an arbitrary order is now placed on \(R^+\), the morphism
\[
 \prod_{\alpha\in R^+}U_\alpha\longrightarrow U
\]
defined by multiplication in \(U\) is still an isomorphism. Indeed, both
sides are flat \(\mathbb Z\)-schemes of finite presentation, so it is
enough to verify the assertion on geometric fibers. We are then reduced
to Lazard's theory (\textit{Bible}, \S13.1, \SGARefLink{sga3:xxv:paragraph:2:1}{Proposition 1}): regard \(U\)
as a group with operators \(T\), and use the fact that the \(U_\alpha\)
are pairwise nonisomorphic as groups with operators, since the characters
\(\alpha\in R^+\) of \(T\) are pairwise distinct on every fiber.


% END inputs/XXV/tex/components/01_section2_1_through_2_5.tex

% BEGIN inputs/XXV/tex/components/02_section2_6_through_2_9.tex
\SGANextSectionTarget{sga3:xxv:subsection:2-6}\SGASubsection{2.6}{}
\SGARefTarget{sga3:xxv:paragraph:2:6}

Replacing \(R^+\) by \(R^-=-R^+\), one similarly constructs the groups
\(U^-\), \(B^-\), and \(U_\alpha\) for \(\alpha\in R^-\), together with
isomorphisms
\[
 p_\alpha\colon\mathbb G_a\xrightarrow{\ \sim\ }U_\alpha,
 \qquad \alpha\in R^-.
\]
Finally introduce the product scheme
\[
 \Omega=U^-\times T\times U.
\]
There is a canonical isomorphism of \(k\)-schemes
\[
 \Omega\otimes_{\mathbb Z}k
 \xrightarrow{\ \sim\ }
 U^-_k\times_kT_k\times_kU_k
 \simeq\Omega_k,
\]
where \(\Omega_k\) is the ``big cell'' of \(G_k\) (\SGARefLink{sga3:xxii:prop:4.1.2}{Exposé XXII, 4.1.2}).

From now on we identify \(\Omega\otimes_{\mathbb Z}k\) with \(\Omega_k\)
by this isomorphism, and regard \(U^-\), \(T\), and \(U\) as subschemes
of \(\Omega\) through their unit sections. Write
\[
 e=((0),e,(0))
\]
for the ``unit section'' of \(\Omega\). Our aim is now to put a
piece-of-group law on \(\Omega\).

\begin{sgastatement}{\SGAInlineRefTarget{sga3:xxv:lemma:2-7}Lemma 2.7}
Let \(\alpha\in\Delta\), and let \(w_{\alpha,k}\) be the element of
\(\operatorname{Norm}_{G_k}(T_k)(k)\) defined by \(X_{\alpha,k}\); recall
that, by definition,
\[
 w_{\alpha,k}
 =p_{-\alpha,k}(-1)p_{\alpha,k}(1)p_{-\alpha,k}(-1).
\]
There exist an open subset \(V_\alpha\) of \(\Omega\), containing the
section \(e\), and a morphism
\[
 h_\alpha\colon V_\alpha\longrightarrow\Omega
\]
such that:
\begin{enumerate}[label=\textup{(\roman*)},leftmargin=2.2em]
\item \(h_\alpha(e)=e\);
\item \((h_\alpha)\otimes_{\mathbb Z}k\) coincides with the restriction
of \(\operatorname{int}(w_{\alpha,k})\) to
\(V_\alpha\otimes_{\mathbb Z}k\subset G_k\);
\item \(T\subset V_\alpha\), and \(h_\alpha\) sends \(T\) into \(T\).
For every \(\beta\in R\), one has \(U_\beta\subset V_\alpha\), and
\(h_\alpha\) sends \(U_\beta\) into \(U_{s_\alpha(\beta)}\).
\end{enumerate}
\end{sgastatement}
\SGARefTarget{sga3:xxv:lemma:2:7}

By the definition of a Chevalley system (\SGARefLink{sga3:XXIII:6.1}{Exposé XXIII, 6.1}), for every
\(\beta\in R\) there exists an integer \(e_\beta=\pm1\) such that
\[
 \operatorname{int}(w_{\alpha,k})p_{\beta,k}(x)
 =p_{s_\alpha(\beta),k}(e_\beta x)
\]
for every \(x\in\mathbb G_a(S)\) and every \(S\to\operatorname{Spec}(k)\).

Let \(S\) be any scheme. Write an arbitrary element of \(\Omega(S)\) in
the form
\[
\left(
 \left(\prod_{\substack{\beta\in R^-\\\beta\ne-\alpha}}
 p_\beta(x_\beta)\right)p_{-\alpha}(x_{-\alpha}),
 \ t,\
 p_\alpha(x_\alpha)
 \left(\prod_{\substack{\beta\in R^+\\\beta\ne\alpha}}
 p_\beta(x_\beta)\right)
\right),
\]
where arbitrary orders have been chosen on \(R^-\setminus\{-\alpha\}\)
and \(R^+\setminus\{\alpha\}\), as permitted by \SGARefLink{sga3:xxv:subsection:2-5}{2.5}. Define a morphism
\(d\colon\Omega\to\mathbb A^1_{\mathbb Z}\) by
\[
 d(u)=\alpha(t)+x_\alpha x_{-\alpha}.
\]

Let \(V_\alpha=\Omega_d\), the open subset of \(\Omega\) defined by the
invertibility of \(d(u)\). It contains \(e\), \(T\), and every \(U_\beta\),
\(\beta\in R\). Define
\[
 h_\alpha\colon V_\alpha\longrightarrow\Omega,\qquad
 h_\alpha(u)=(a(u),b(u),c(u)),
\]
where
\begin{align*}
a(u)
 &=
\left(\prod_{\substack{\beta\in R^-\\\beta\ne-\alpha}}
 p_{s_\alpha(\beta)}(e_\beta x_\beta)\right)
p_{-\alpha}\!\left(-x_\alpha d(u)^{-1}\right),\\
b(u)&=t\cdot\alpha^\ast(d(u)),\\
c(u)
 &=p_\alpha\!\left(-x_{-\alpha}d(u)^{-1}\right)
\left(\prod_{\substack{\beta\in R^+\\\beta\ne\alpha}}
 p_{s_\alpha(\beta)}(e_\beta x_\beta)\right).
\end{align*}
Since \(s_\alpha\) permutes the positive roots distinct from \(\alpha\)
(respectively the negative roots distinct from \(-\alpha\)), \(c(u)\)
(respectively \(a(u)\)) is a section of \(U\) (respectively \(U^-\)),
so the preceding morphism is well defined. It plainly satisfies (i)
and (iii). Assertion (ii) follows immediately from the definition of
the \(e_\beta\), \(\beta\in R\), and \SGARefLink{sga3:xx:3.12}{Exposé XX, 3.12}.

\begin{sgastatement}{\SGAInlineRefTarget{sga3:xxv:lemma:2-8}Lemma 2.8}
There exist open subsets \(V,V'\) of \(\Omega\) and morphisms
\[
 h\colon V\longrightarrow\Omega,\qquad
 h'\colon V'\longrightarrow\Omega
\]
with the following properties:
\begin{enumerate}[label=\textup{(\roman*)},leftmargin=2.2em]
\item \(V\) and \(V'\) contain \(e\), and \(h(e)=h'(e)=e\);
\item\SGAInlineRefTarget{sga3:xxv:lemma:2-8:item:ii} the morphism induced by
\[
 h'\circ h\colon h^{-1}(V')\longrightarrow\Omega
\]
is the restriction of the identity morphism;
\item\SGAInlineRefTarget{sga3:xxv:lemma:2-8:item:iii} the \(k\)-morphisms \(h\otimes_{\mathbb Z}k\) and
\(h'\otimes_{\mathbb Z}k\) are the restrictions, to
\(V\otimes_{\mathbb Z}k\) and \(V'\otimes_{\mathbb Z}k\), of
automorphisms of the group \(G_k\);
\item\SGAInlineRefTarget{sga3:xxv:lemma:2-8:item:iv} \(V\) and \(V'\) contain \(U,T,U^-\); the morphisms \(h,h'\)
send \(U\) into \(U^-\), \(U^-\) into \(U\), and \(T\) into \(T\).
\end{enumerate}
\end{sgastatement}
\SGARefTarget{sga3:xxv:lemma:2:8}

Let \(w^0\) be the element of the Weyl group of \(G_k\) that sends \(R^+\)
to \(R^-\). Write
\[
 w^0=s_{\alpha_n}\cdots s_{\alpha_1},\qquad \alpha_i\in\Delta
\]
(with no relation to the numbering in \SGARefLink{sga3:xxv:subsection:2-1}{2.1}), and put
\[
 w_0=w_{\alpha_n,k}\cdots w_{\alpha_1,k}
 \in\operatorname{Norm}_{G_k}(T_k)(k).
\]
Define recursively, for \(i\leq n\), an open subset \(V_i\) of \(\Omega\)
and a morphism \(h_i\colon V_i\to\Omega\) by
\[
 V_0=\Omega,\qquad h_0=\operatorname{id},
\]
and, for \(i=0,\ldots,n-1\),
\[
 V_{i+1}=h_i^{-1}(V_{\alpha_{i+1}}),\qquad
 h_{i+1}=h_{\alpha_{i+1}}\circ h_i,
\]
where \(V_{\alpha_j}\) and \(h_{\alpha_j}\) are those of \SGARefLink{sga3:xxv:lemma:2-7}{2.7}.

Take \(V=V_n\) and \(h=h_n\). The conditions (i), (iii), and (iv) concerning
\(V\) and \(h\) hold: for (i) and (iii) this follows immediately from
\SGARefLink{sga3:xxv:lemma:2-8}{2.8}, and for (iv) from the fact that \(h\otimes_{\mathbb Z}k\) is the
restriction of \(\operatorname{int}(w_0)\) to
\(V\otimes_{\mathbb Z}k\).

Since \((w^0)^2=1\), we also have
\[
 w^0=s_{\alpha_1}\cdots s_{\alpha_n}
 =(s_{\alpha_1}^{\,3})\cdots(s_{\alpha_n}^{\,3}).
\]
Put
\[
 w'_0=(w_{\alpha_1,k})^3\cdots(w_{\alpha_n,k})^3.
\]
Repeating the construction above gives an open subset \(V'\) and a
morphism \(h'\) that also satisfy (i), (iii), and (iv). Moreover,
\(h'\otimes_{\mathbb Z}k\) is the restriction of
\(\operatorname{int}(w'_0)\) to \(V'\otimes_{\mathbb Z}k\). For every
simple root \(\alpha\in\Delta\), however,
\((w_{\alpha,k})^4=e\) (\SGARefLink{sga3:xx:3.1}{Exposé XX, 3.1}), hence \(w'_0w_0=e\).
It follows that \(h'\circ h\) induces the identity morphism on a nonempty
open subset of \(\Omega\otimes_{\mathbb Z}k\). Since \(\Omega\) is smooth
and of finite presentation over \(\mathbb Z\),
\(\Omega\otimes_{\mathbb Z}k\) is schematically dense in \(\Omega\);
this proves (ii).

\begin{sgastatement}{\SGAInlineRefTarget{sga3:xxv:proposition:2-9}Proposition 2.9}
There exist an open subset \(V_1\) of \(\Omega\times\Omega\), an open
subset \(V_2\) of \(\Omega\), and morphisms
\[
 \pi\colon V_1\longrightarrow\Omega,\qquad
 \sigma\colon V_2\longrightarrow\Omega
\]
having the following properties:
\begin{enumerate}[label=\textup{(\roman*)},leftmargin=2.2em]
\item if \(x\in\Omega(S)\), then \((e,x)\) and \((x,e)\) are sections
of \(V_1\), and
\[
 \pi(e,x)=\pi(x,e)=x;
\]
\item \(V_2\) contains \(e\), and \(\sigma(e)=e\);
\item\SGAInlineRefTarget{sga3:xxv:proposition:2-9:item:iii} \(\pi_k\) and \(\sigma_k\) are the restrictions of the morphisms
\(G_k\times_kG_k\to G_k\) and \(G_k\to G_k\) defined respectively by
multiplication and inversion.
\end{enumerate}
\end{sgastatement}
\SGARefTarget{sga3:xxv:proposition:2:9}

\SGAProof
Let \((v,t,u)\) and \((v',t',u')\) be two sections of \(\Omega\).
By \SGARefLink{sga3:xxv:lemma:2-8:item:iv}{2.8(iv)}, \(h(u)\) is a section of \(U^-\) and \(h(v')\) a section
of \(U\); thus \((h(u),e,h(v'))\) may be regarded as a section of
\(\Omega\). Define \(V_1\) as the open subset of \(\Omega\times\Omega\)
specified by
\[
\bigl((v,t,u),(v',t',u')\bigr)\in V_1(S)
\quad\Longleftrightarrow\quad
(h(u),e,h(v'))\in V'(S),
\]
using the notation of \SGARefLink{sga3:xxv:lemma:2-8}{2.8}. If the pair on the left is a section of
\(V_1\), then \(h'(h(u),e,h(v'))\) is defined and is a section of
\(\Omega\); write
\[
 h'(h(u),e,h(v'))=(v'',t'',u'').
\]
Set
\[
 \pi\bigl((v,t,u),(v',t',u')\bigr)
 =
 \bigl(v\cdot t v''t^{-1},\ tt''t',\
 t'^{-1}u''t'\cdot u'\bigr).
\]
Assertion (i) follows immediately from \SGARefLink{sga3:xxv:lemma:2-8:item:ii}{2.8(ii)}. To verify the assertion
about \(\pi\) in (iii), observe that
\[
 h'(h(u),e,h(v'))=uv'=v''t''u''
\]
whenever \(u\in U(S)\), \(v'\in U^-(S)\), and \(S\to k\), by
\SGARefLink{sga3:xxv:lemma:2-8:item:iii}{2.8(iii)} and (ii).

The morphism \(\sigma\) is constructed similarly. If \((v,t,u)\) is a
section of \(\Omega\), then \(h(u^{-1})\) is a section of \(U^-\),
\(h(v^{-1})\) a section of \(U\), and \(h(t^{-1})\) a section of \(T\).
Thus
\[
 (h(u^{-1}),h(t^{-1}),h(v^{-1}))
\]
is a section of \(\Omega\).
Define an open subset \(V_2\) of \(\Omega\) by
\[
 (v,t,u)\in V_2(S)
 \quad\Longleftrightarrow\quad
 (h(u^{-1}),h(t^{-1}),h(v^{-1}))\in V'(S),
\]
and define
\[
 \sigma(v,t,u)=h'(h(u^{-1}),h(t^{-1}),h(v^{-1})).
\]
The required properties of \(\sigma\) are verified in the same way.

% END inputs/XXV/tex/components/02_section2_6_through_2_9.tex

% BEGIN inputs/XXV/tex/components/03_section2_10_and_section3.tex
\begin{sgastatement}{Corollary 2.10}
\(\pi\) is ``generically associative'' and \(\sigma\) is a ``generic
inverse'': if \(x,y,z\in\Omega(S)\), and if the expressions below are
defined (as they always are over an open subset of \(\Omega\) containing
the unit section), then
\[
 \pi(x,\pi(y,z))=\pi(\pi(x,y),z),\qquad
 \pi(x,\sigma(x))=e=\pi(\sigma(x),x).
\]
\end{sgastatement}
\SGARefTarget{sga3:xxv:corollary:2:10}

Indeed, the two sides of each formula define morphisms between smooth
\(\mathbb Z\)-schemes of finite presentation, and the morphisms agree on
the generic fibers by \SGARefLink{sga3:xxv:proposition:2-9:item:iii}{2.9(iii)}.\footnote{The French source reads
``by \SGARefLink{sga3:xxv:proposition:2-9:item:iii}{2.10(iii)}''; the intended reference is Proposition 2.9(iii).}

\begin{sgastatement}{\SGAInlineRefTarget{sga3:xxv:corollary:2-11}Corollary 2.11}
Let \(\alpha\in R\), and let \(S\) be a scheme. Suppose that
\(x,y\in\mathbb G_a(S)\) satisfy
\[
 (p_\alpha(x),p_{-\alpha}(y))\in V_1(S),
 \qquad 1+xy\in\mathbb G_m(S).
\]
These conditions define an open subset of \(\mathbb G_{a,S}^2\)
containing the section \((0,0)\). Then
\[
 \pi(p_\alpha(x),p_{-\alpha}(y))
 =
 \left(
 p_{-\alpha}\!\left(\frac{y}{1+xy}\right),
 \alpha^\ast(1+xy),
 p_\alpha\!\left(\frac{x}{1+xy}\right)
 \right).
\]
\end{sgastatement}
\SGARefTarget{sga3:xxv:corollary:2:11}

The proof is the same as above, by \SGARefLink{sga3:xx:2.1}{Exposé XX, 2.1}.

\SGASection{3}{Existence theorem: conclusion of the proof}
\SGARefTarget{sga3:xxv:section:3}

For convenience, make the following definition.

\begin{sgastatement}{\SGAInlineRefTarget{sga3:xxv:definition:3-1}Definition 3.1}
Let \(S\) be a scheme and \(G\) an \(S\)-group scheme. The group \(G\)
is called \emph{admissible} if there is an open immersion of
\(S\)-schemes
\[
 i\colon\Omega_S=\Omega\times S\longrightarrow G
\]
such that:
\begin{enumerate}[label=\textup{(\roman*)},leftmargin=2.2em]
\item the diagram
\[
\begin{tikzcd}[column sep=4.5em,row sep=3.2em]
(V_1)_S
  \arrow[r]
  \arrow[d,"\pi"]
&
\Omega_S\times_S\Omega_S
  \arrow[r,"i\times_S i"]
&
G\times_SG
  \arrow[d,"\pi_G"]
\\
\Omega_S
  \arrow[rr,"i_S"]
&&
G
\end{tikzcd}
\]
is commutative, where \(\pi_G\) denotes the multiplication morphism
of \(G\);
\item there is a finite set of sections \(a_j\in\Omega(S)\) such that
the open subsets \(i(a_j)\,i(\Omega_S)\) cover \(G\).
\end{enumerate}
\end{sgastatement}
\SGARefTarget{sga3:xxv:definition:3:1}

By ``Weil's theorem'' (\SGARefLink{sga3:xviii:corollary:3-13:item:iii}{Exposé XVIII, 3.13(iii)} and (iv)), we have:

\begin{sgastatement}{Lemma 3.2}
If every admissible \(S\)-group is affine whenever \(S\) is étale and
of finite type over \(\mathbb Z\), then there exists an admissible
affine \(\mathbb Z\)-group.
\end{sgastatement}
\SGARefTarget{sga3:xxv:lemma:3:2}

We now have:

\begin{sgastatement}{Lemma 3.3}
Let \(S\) be a scheme and \(G\) an admissible \(S\)-group. Then \(G\)
is smooth and of finite presentation over \(S\), with affine,
connected, semisimple fibers.
\end{sgastatement}
\SGARefTarget{sga3:xxv:lemma:3:3}

Since \(\Omega_S\) is smooth and of finite presentation over \(S\),
with connected fibers, the same is true of \(G\) by condition (ii).
To prove \SGARefLink{sga3:xxv:lemma:3:3}{3.3}, we may therefore suppose that \(S\) is the spectrum of a
field \(K\). Identify \(\Omega_K\) with its image in \(G\). Plainly
\(\Omega_K\) is the product
\[
 \left(\prod_{\alpha\in R^-}U_{\alpha,K}\right)
 T_K
 \left(\prod_{\alpha\in R^+}U_{\alpha,K}\right)
\]
of the subgroups \(T_K\) and \(U_{\alpha,K}\), \(\alpha\in R\), of
\(G\). Thus the Lie algebra of \(G\) is identified with the direct sum
\[
 \operatorname{Lie}(T_K)
 \oplus
 \bigoplus_{\alpha\in R}\operatorname{Lie}(U_{\alpha,K}).
\]
The inner automorphism defined by a section of \(T_K\) acts on
\(U_{\alpha,K}\), and hence on \(\operatorname{Lie}(U_{\alpha,K})\),
through the character
\[
 \alpha\in R\subset M
 \simeq\operatorname{Hom}_{K\text{-gr.}}(T_K,\mathbb G_{m,K}).
\]
The preceding decomposition of \(\operatorname{Lie}(G)\) is therefore
exactly its decomposition under the adjoint action of \(T\). Thus the
roots of \(G_K\) with respect to \(T_K\) are the elements
\(\alpha\in R\).

Apply \SGARefLink{sga3:xix:corollary:1-13}{Exposé XIX, 1.13}. Let \(T_\alpha\) be the maximal torus of
\(\ker(\alpha)\subset T_K\), and put
\[
 Z_\alpha=\operatorname{Centr}_G(T_\alpha).
\]
It is enough to prove that each \(Z_\alpha\) is reductive. But
\[
 Z_\alpha\cap\Omega_K
 =
 \left(
 \prod_{\substack{\beta\in R^-\\\beta|_{T_\alpha}=e}}
 U_{\beta,K}
 \right)
 T_K
 \left(
 \prod_{\substack{\beta\in R^+\\\beta|_{T_\alpha}=e}}
 U_{\beta,K}
 \right).
\]
The roots vanishing on \(T_\alpha\) are the rational multiples of
\(\alpha\), hence precisely \(\alpha\) and \(-\alpha\). Consequently,
\[
 Z_\alpha\cap\Omega_K
 =U_{-\alpha,K}\,T_K\,U_{\alpha,K}.
\]
By \SGARefLink{sga3:xx:3.4}{Exposé XX, 3.4}, to prove \(Z_\alpha\) reductive it is enough to show
that \(U_{\alpha,K}\) and \(U_{-\alpha,K}\) do not commute, and this
follows immediately from \SGARefLink{sga3:xxv:corollary:2-11}{2.11}.

It follows from \SGARefLink{sga3:xxv:lemma:3:2}{3.2} and \SGARefLink{sga3:xxv:lemma:3:3}{3.3} that the proof will be complete once the
following lemma is proved.

\begin{sgastatement}{Lemma 3.4}
Let \(S\) be a locally noetherian scheme of dimension at most \(1\),
and let \(G\) be a smooth \(S\)-group of finite type, with affine,
connected, semisimple fibers. Then \(G\) is affine (and hence
semisimple).
\end{sgastatement}
\SGARefTarget{sga3:xxv:lemma:3:4}

\noindent\textit{Nota.}\quad
In Exposé XVI it was shown that \SGARefLink{sga3:xxv:lemma:3:4}{3.4} holds without any hypothesis on
\(S\), but the proof is comparatively delicate. Since only the special
case \SGARefLink{sga3:xxv:lemma:3:4}{3.4} is needed here, we give a direct proof.

Consider the Lie algebra \(\mathfrak g\) of \(G\), a locally free
\(\mathcal O_S\)-module, and the adjoint representation
\[
 \operatorname{Ad}\colon
 G\longrightarrow\operatorname{GL}_{\mathcal O_S}(\mathfrak g).
\]
To prove that \(G\) is affine over \(S\), it is enough to prove that
\(\operatorname{Ad}\) is affine. Since \(G\) is smooth with connected
fibers, it is separated over \(S\) (Exposé VI\textsubscript{B}, \SGARefLink{cor:sga3-VIB-5.5}{5.5}),
so \(\operatorname{Ad}\) is separated. By the result proved in the
Appendix (see \SGARefLink{sga3:xxv:proposition:4:1}{4.1}), it suffices to show that \(\operatorname{Ad}\) is
quasi-finite. We are thus reduced to the case in which \(S\) is the
spectrum of a field. In that case \(G\) is affine and therefore
semisimple, and the assertion is \SGARefLink{sga3:xxii:prop:5.7.14}{Exposé XXII, 5.7.14}.

% END inputs/XXV/tex/components/03_section2_10_and_section3.tex

% BEGIN inputs/XXV/tex/components/04_section4_and_bibliography.tex
\SGASection{4}{Appendix}
\SGARefTarget{sga3:xxv:section:4}

We used the following proposition in the course of the proof.

\begin{sgastatement}{Proposition 4.1}
Let \(S\) be a locally noetherian scheme of dimension at most \(1\);
let \(G\) and \(H\) be two \(S\)-group schemes of finite type; and let
\(f\colon G\to H\) be a quasi-finite separated group homomorphism.
If \(G\) is flat over \(S\),\footnote{Editors' note: We have added the
flatness hypothesis, which had been omitted.} then \(f\) is affine.
\end{sgastatement}
\SGARefTarget{sga3:xxv:proposition:4:1}

We shall give the proof only when \(G\) is smooth over \(S\), which is
indeed the hypothesis satisfied in the application above.

\SGANextSectionTarget{sga3:xxv:subsection:4-2}\SGASubsection{4.2}{}
\SGARefTarget{sga3:xxv:paragraph:4:2}

By EGA II, 1.6.4, we may suppose that \(S\) is reduced. By the usual
techniques of passage to the limit,\footnote{Editors' note: See
EGA IV\textsubscript{3}, 8.10.5(viii).} we may suppose \(S\) local. If
\(\dim(S)=0\), the assertion is trivial.\footnote{Editors' note:
Indeed, if \(S=\operatorname{Spec}(k)\), with \(k\) a field, then \(f\)
is the composite of the projection \(p\colon G\to G/\ker(f)\) and a
closed immersion \(i\). Since \(\ker(f)\) is finite over \(k\), \(p\)
is finite (Exposé VI\textsubscript{B}, 9.2).}
Suppose therefore that \(\dim(S)=1\). By faithfully flat descent, we may
suppose that \(S\) is complete with algebraically closed residue field.
After replacing \(S\) by its normalization \(\widetilde S\), if
necessary, we may suppose \(S\) normal (EGA II, 6.7.1, and
EGA \(0_{\mathrm{IV}}\), 23.1.5).\footnote{Editors' note:
Indeed, the irreducible components \(S_1,\ldots,S_r\) of \(S\) are
complete integral noetherian local schemes. By a theorem of Nagata
(EGA \(0_{\mathrm{IV}}\), 23.1.5), the normalization
\(\widetilde S_i\) is finite over \(S_i\), and hence
\(\widetilde S\) is finite over \(S\). A theorem of Chevalley
(EGA II, 6.7.1) then shows that if
\(f\times_S\widetilde S\) is affine, so is \(f\).}
We are thus reduced to the case in which \(S\) is the spectrum of a
complete discrete valuation ring \(A\) with algebraically closed
residue field.

\SGANextSectionTarget{sga3:xxv:subsection:4-3}\SGASubsection{4.3}{}
\SGARefTarget{sga3:xxv:paragraph:4:3}

Let \(\eta\) (respectively \(s\)) be the generic (respectively closed)
point of \(S\). Consider the image \(f_\eta(G_\eta)\) of \(G_\eta\)
in \(H_\eta\). It is a closed subgroup scheme of \(H_\eta\). Let
\(H'\) be the schematic closure of \(f_\eta(G_\eta)\) in \(H\).
Since \(H'\to H\) is affine (it is a closed immersion), we may replace
\(H\) by \(H'\), and thus suppose that \(H\) is flat over \(S\) and
that \(f_\eta\) is surjective. Since \(f_\eta\) is finite and \(G,H\)
are flat over \(S\), we have
\[
 \dim(G_s)=\dim(G_\eta)=\dim(H_\eta)=\dim(H_s).
\]

\SGANextSectionTarget{sga3:xxv:subsection:4-4}\SGASubsection{4.4}{}
\SGARefTarget{sga3:xxv:paragraph:4:4}

Let \(H_s^0,\ldots,H_s^n\) be the irreducible components of \(H_s\),
where \(H_s^0\) denotes the identity component, and let
\(z_0,\ldots,z_n\) be their generic points. Since each local ring
\(\mathcal O_{H,z_i}\) has dimension at most \(1\), the morphism
\[
 G\times_H\mathcal O_{H,z_i}\longrightarrow\mathcal O_{H,z_i}
\]
is affine, being quasi-finite and separated (see \SGARefLink{sga3:xvi:lemma:4:2}{Exposé XVI, Lemma 4.2}).
Thus \(G\) is affine over \(H\) in a neighborhood of \(z_i\). Let \(V\)
be the largest open subset of \(H\) over which \(G|_V\) is affine.
Then \(V\) contains every \(z_i\),\footnote{Editors' note: The following
sentence has been added.} and therefore contains at least one closed
point \(y_i\in H_s^i\). Moreover,
\(\kappa(y_i)=\kappa(s)\), since \(\kappa(s)\) is algebraically closed.

On the other hand, \(V\) is plainly stable under translation by any
element \(g\in G(S)\). We have \(\dim(G_s)=\dim(H_s)\), and \(f_s\) is
finite, so
\[
 f(s)\colon G_s^0(s)\longrightarrow H_s^0(s)
\]
is surjective. Since \(A\) is complete and \(G\) is smooth over \(S\),
the canonical map \(G^0(S)\to G_s^0(s)\) is surjective. Since
\(H_s^0(s)\) acts transitively on each \(H_s^i(s)\), it follows that
\(V\supset H_s(s)\), and therefore, since \(\kappa(s)\) is algebraically
closed, \(V\supset H_s\).\footnote{Editors' note: When \(G\) is not
assumed smooth over \(S\), one may argue as follows. Let \(h_0\) be a
rational point of \(H_s^0\), and let \(g_0\) be a rational point of
\(G_s^0\) such that \(f(g_0)=h_0\). By Exposé VI\textsubscript{B},
5.6.1, there is a commutative diagram
\[
\begin{tikzcd}[ampersand replacement=\&,column sep=4.0em,row sep=3.0em]
S'' \arrow[r,"g"] \arrow[d,"\pi"'] \arrow[dr,"\phi"]
  \& G \arrow[d]\\
S' \arrow[r,"w"] \& S
\end{tikzcd}
\]
in which \(w\) is étale and surjective, \(\pi\) is finite and
surjective, and \(\phi^{-1}(s)\) consists of a single point \(s''\)
such that \(g(s'')=g_0\). Then \(G_{S''}\to H_{S''}\) is affine above
a neighborhood of \(h_0y_i\), and the same is true of
\(G_{S'}\to H_{S'}\) (EGA II, 6.7.1), and hence of \(G\to H\) by
faithfully flat descent (EGA IV\textsubscript{2}, 2.7.1(xiii)).
Thus \(h_0y_i\in V\), from which it follows that \(V\) contains
\(H_s(s)\), and therefore \(H_s\).}

Clearly \(V\supset H_\eta\), since \(f_\eta\) is finite. Consequently,
\(V=H\). \SGAQED

\section*{Bibliography\footnote{Editors' note: Additional references
cited in this Exposé.}}
\addcontentsline{toc}{section}{Bibliography}
\SGARefTarget{sga3:xxv:bibliography}

\begin{description}[style=multiline,leftmargin=4.8em,labelwidth=4.0em]
\item[\textup{[BLie]}]
N.~Bourbaki, \textit{Groupes et algèbres de Lie}, Chapters I and
VII--VIII, Hermann, 1960 and 1975.

\item[\textup{[BT84]}]
F.~Bruhat and J.~Tits, \textit{Groupes réductifs sur un corps local II.
Schémas en groupes. Existence d'une donnée radicielle valuée},
Publ.\ Math.\ I.H.É.S. \textbf{60} (1984), 5--184.

\item[\textup{[DG70]}]
M.~Demazure and P.~Gabriel, \textit{Groupes algébriques},
Masson \& North-Holland, 1970.

\item[\textup{[Ko66]}]
B.~Kostant, \textit{Groups over \(\mathbb Z\)}, pp.~90--98 in
\textit{Algebraic groups and their discontinuous subgroups}
(A.~Borel and G.~D.~Mostow, eds.), Proc.\ Symp.\ Pure Math.\
\textbf{IX}, American Mathematical Society, 1966.

\item[\textup{[Lu09]}]
G.~Lusztig, \textit{Study of a \(\mathbb Z\)-form of the coordinate
ring of a reductive group}, J.\ Amer.\ Math.\ Soc.\ \textbf{22}
(2009), no.~3, 739--769.

\item[\textup{[Se66]}]
J.-P.~Serre, \textit{Algèbres de Lie semi-simples complexes},
Benjamin, 1966.

\item[\textup{[Sp98]}]
T.~A.~Springer, \textit{Linear algebraic groups}, 2nd ed.,
Birkhäuser, 1998.

\item[\textup{[Ti66]}]
J.~Tits, \textit{Sur les constantes de structure et le théorème
d'existence des algèbres de Lie semi-simples},
Publ.\ Math.\ I.H.É.S. \textbf{31} (1966), 21--58.
\end{description}

% END inputs/XXV/tex/components/04_section4_and_bibliography.tex


